\documentclass[11pt]{article}
\usepackage[margin=1in]{geometry}
\usepackage{amsmath,amssymb,amsthm,enumitem,hyperref}
\hypersetup{colorlinks=true,urlcolor=blue,linkcolor=blue}
\newtheorem{theorem}{Theorem}
\title{Increasing minima of the normalized Bessel Tur\'anian}
\author{Principia Math In Progress}
\date{21 September 2026}
\begin{document}
\maketitle
\noindent\textbf{Task 1.} Resolve Baricz's minimum-value conjecture
\cite[Conjecture 4.2]{B}, repeated in \cite[Conjecture 2.3]{BP} and recorded as
\href{https://mathdb.com/p/321003}{MathDB 321003}.
The database attributes the conjecture to B\'ar\'any; the primary sources
attribute it to Baricz. The printed strict bound must allow equality at its
defining minimizer. We prove the intended sharp bound and strict increase
of all its constants; neither transcription issue is the mathematical result.
Review is by the authoring agent only.

Fix $\nu>0$, let $j_{\nu,n}$ be the positive zeros of $J_\nu$, and put
\[
 \Phi_\nu(x)=1-\frac{J_{\nu-1}(x)J_{\nu+1}(x)}{J_\nu(x)^2},
 \qquad x>0,\quad J_\nu(x)\ne0.
\]
\begin{theorem}
Each interval $(j_{\nu,n},j_{\nu,n+1})$, $n\ge1$, contains exactly one
critical point $\alpha_{\nu,n}$ of $\Phi_\nu$, its strict global minimum on
that interval. With $\beta_{\nu,n}=\Phi_\nu(\alpha_{\nu,n})$ and
$\beta_{\nu,0}=1/(\nu+1)$, one has
\[
 \frac1{\nu+1}<\beta_{\nu,1}<\beta_{\nu,2}<\cdots<1,
 \qquad \lim_{n\to\infty}\beta_{\nu,n}=1.
\]
For $x\in(j_{\nu,n},j_{\nu,n+1})$,
\[
 J_\nu(x)^2-J_{\nu-1}(x)J_{\nu+1}(x)
 \ge \beta_{\nu,n}J_\nu(x)^2,
\]
with equality exactly at $x=\alpha_{\nu,n}$. On $(0,j_{\nu,1})$ the
corresponding bound with $\beta_{\nu,0}$ is strict.
\end{theorem}

\paragraph{Existence and uniqueness, using the known zero sum.}
The classical identity recorded in \cite[Section 4.2]{B} and
\cite[Section 2.8]{BP} is
\begin{equation}\label{sum}
 \Phi_\nu(x)=4\sum_{k\ge1}\frac{j_{\nu,k}^2}{(x^2-j_{\nu,k}^2)^2}.
\end{equation}
Baricz already proves strict convexity in $x$ on every zero gap in
\cite[p.~262]{B}; thus the convexity ingredient is prior work.
For completeness, write $a_k=j_{\nu,k}^2$ and $H(t)=\Phi_\nu(\sqrt t)$.
Then
\begin{equation}\label{convex}
 H'(t)=8\sum_{k\ge1}\frac{a_k}{(a_k-t)^3},\qquad
 H''(t)=24\sum_{k\ge1}\frac{a_k}{(a_k-t)^4}>0.
\end{equation}
The standard zero asymptotic $j_{\nu,k}\sim\pi k$ makes the series and these
derivatives locally uniformly convergent away from poles: their tails are
respectively $O(k^{-2})$, $O(k^{-4})$, and $O(k^{-6})$ termwise.

\newpage
On $(a_n,a_{n+1})$, $H'$ tends to $-\infty$ at the left endpoint and
$+\infty$ at the right endpoint. The other terms are bounded near either
fixed endpoint. Hence \eqref{convex} gives exactly one zero, with $H'$
negative before it and positive afterward. Since $x\mapsto x^2$ is
strictly increasing for $x>0$, this proves the asserted minimum and its
equality case. On $(0,a_1)$ all terms in $H'$ are positive. The Bessel
series at zero, or the Rayleigh sum in \cite{B}, gives
$\Phi_\nu(0+)=1/(\nu+1)$, proving the first-interval assertion.

\paragraph{A scalar relation orders the minimum values.}
On any zero-free interval define
\[
 r(x)=\frac{xJ_\nu'(x)}{J_\nu(x)}.
\]
The usual adjacent-order recurrences and Bessel differential equation give
\begin{equation}\label{riccati}
 \Phi_\nu(x)=1+\frac{r(x)^2-\nu^2}{x^2},\qquad
 r'(x)=\frac{\nu^2-x^2-r(x)^2}{x}=-x\Phi_\nu(x).
\end{equation}
These are classical identities; compare \cite[Eqs.~(17) and following]{BPS}.
Differentiating the first identity and substituting the second yields
\begin{equation}\label{stationary}
 \Phi_\nu'(x)=-\frac2x\bigl((1+r(x))\Phi_\nu(x)-1\bigr).
\end{equation}
Fix $n\ge1$ and abbreviate $\alpha=\alpha_{\nu,n}$, $b=\beta_{\nu,n}$,
and $s=r(\alpha)$. By \eqref{sum}, $b>0$. Rolle's theorem supplies a
$c\in(j_{\nu,n},j_{\nu,n+1})$ with $J_\nu'(c)=0$. Since $J_\nu(c)\ne0$,
\eqref{riccati} gives $\Phi_\nu(c)=1-\nu^2/c^2<1$. Therefore $b<1$.
At the minimum, \eqref{stationary} now implies
\[
 b(1+s)=1,\qquad s=b^{-1}-1>0.
\]
Substituting $b=1/(1+s)$ into \eqref{riccati} gives
\begin{equation}\label{curve}
 \alpha^2=F_\nu(s):=\frac{(\nu^2-s^2)(1+s)}s
 =\nu^2+\frac{\nu^2}s-s-s^2.
\end{equation}
Positivity of $\alpha^2$ implies $s<\nu$. Moreover,
\begin{equation}\label{slope}
 F_\nu'(s)=-\frac{\nu^2}{s^2}-1-2s<0\qquad(s>0).
\end{equation}
The points $\alpha_{\nu,n}$ strictly increase with $n$, because they lie
in consecutive disjoint zero gaps. Equations \eqref{curve}--\eqref{slope}
therefore force $s_n=r(\alpha_{\nu,n})$ to strictly decrease, so
$\beta_{\nu,n}=1/(1+s_n)$ strictly increases. The bound $s_n<\nu$ gives
$\beta_{\nu,n}>1/(\nu+1)$. Finally $\alpha_{\nu,n}\to\infty$ and
\eqref{curve} force $s_n\to0$: otherwise the decreasing positive sequence
$s_n$ would have a positive limit at which $F_\nu$ is finite. Thus
$\beta_{\nu,n}\to1$, completing the proof.\hfill$\square$

\newpage
\paragraph{Scope and attribution.}
The complete 16-page current preprint of \cite{BP}, its complete 28-page
typeset publisher author copy (marked ``Ahead of Print'' with the journal
DOI), and all 11 pages of the original journal paper \cite{B} were read.
The later 11-page journal paper \cite{BPS} was also read completely; its
extrema theorem concerns the unnormalized Tur\'anian and does not assert
the present increase of normalized minima. Its logarithmic-derivative
identities are credited above. The zero sum, convexity, first-interval
bound, and the known bound $\beta_{\nu,n}<1$ are prior ingredients.
The contribution of this proof is their completion by the stationary
relation \eqref{curve}, which orders every minimum for every $\nu>0$.

The printed strict inequality at $\alpha_{\nu,n}$ in both conjecture
statements is inconsistent with the definition of $\beta_{\nu,n}$.
The theorem gives the sharp intended version, including the exact
equality set, and settles the substantive monotonicity assertion.
No additional credit is assigned for this correction or for the limit.

\paragraph{Verification.}
The companion script checks the differential identities by exact
substitution of the Bessel equation, the zero-sum derivative and scalar
curve identities symbolically, and their order implications on rational
parameters. Direct high-precision Bessel evaluations independently check
representative minima, the source's four numerical locations at $\nu=2$,
and the derivative signs on both sides. These finite computations support
the written proof; they do not establish its universal quantifiers.
Initial and refreshed MathDB and bounded public-claim checks are archived.
No independent peer review or proof-assistant verification is claimed.

\paragraph{Tasks 2--5 (proposed, not solved here).}
\begin{enumerate}[label=Task \arabic*.,start=2,leftmargin=4.8em]
\item Derive an asymptotic expansion for $\alpha_{\nu,n}$ and
$\beta_{\nu,n}$ as $n\to\infty$, with explicit error bounds uniform for
$\nu$ in a prescribed compact subset of $(0,\infty)$.
\item Determine whether $\nu\mapsto\beta_{\nu,n}$ is strictly decreasing
on $(0,\infty)$ for each fixed $n\ge1$.
\item Classify the corresponding sequences of normalized minima for
$C_\nu=\cos\theta\,J_\nu-\sin\theta\,Y_\nu$, including the first interval
and the admissible parameter range.
\item Obtain explicit two-sided bounds for
$\beta_{\nu,n+1}-\beta_{\nu,n}$ in terms of $\nu$ and adjacent Bessel
zeros, with asymptotically matching leading terms.
\end{enumerate}

\begin{thebibliography}{9}\small
\bibitem{B} \'A. Baricz, Tur\'an type inequalities for modified Bessel
functions, \emph{Bull. Aust. Math. Soc.} \textbf{82} (2010), 254--264.
\href{https://doi.org/10.1017/S000497271000002X}{doi:10.1017/S000497271000002X}.
\bibitem{BP} \'A. Baricz and T. K. Pog\'any, Tur\'an determinants of
Bessel functions, \emph{Forum Math.} \textbf{26} (2014), 295--322.
\href{https://doi.org/10.1515/form.2011.160}{doi:10.1515/form.2011.160};
\href{https://arxiv.org/abs/1101.4624v2}{arXiv:1101.4624v2}.
\bibitem{BPS} \'A. Baricz, S. Ponnusamy, and S. Singh, Tur\'an type
inequalities for general Bessel functions, \emph{Math. Inequal. Appl.}
\textbf{19} (2016), 709--719.
\href{https://doi.org/10.7153/mia-19-51}{doi:10.7153/mia-19-51}.
\end{thebibliography}
\end{document}
